\section{Conclusions}

The project analysed custome-support-tickets data from Hugging Face to identify the main issues raised by users and to evaluate how different text analytics methods influence the insights extracted. The main two axes were investigated: clustering (Axis 1) and topic modelling (Axis 2).

The research ultimately yielded several categories of customer concerns across different methods, though the exact structure of these categories varies depending on the method used, the themes appear consistently across both clustering and topic modelling results. The categories include technical problems, such as system errors and software updates; secure and medical concerns such as access, security, and data management; and broader business topics, such as analytics, strategy, and digital operations; project issues such as integration, management, and data tools. 

In terms of clustering performance, the results show that the choice of text representation has a strong impact on clustering. In particular, Embeddings combined with K-means achieve the highest silhouette score, indicating stronger semantic grouping of customer issues. In contrast, TF-IDF combined with hierarchical clustering produces more stable and consistent cluster structures, as reflected by higher ARI and NMI scores. In practice, embeddings are more useful for uncovering meaningful groupings, while TF-IDF can provide more predictable and interpretable results. For topic modelling, the comparison between two approaches highlights a similar contrast. LDA produces topics that are generally clearer and easier to interpret, with keywords that directly reflect common user issues. BERTopic, on the other hand, aligns more closely with the structure observed in clustering, indicating that it captures relationships in the embedding space more effectively. However, it also produces imbalanced topic distributions in some cases, which reduces interpretability. Overall, the results demonstrate that no single method is optimal in all aspects. Embedding-based approaches are more effective at capturing semantic relationships, while TF-IDF provide more stable structures. Similarly, LDA is better suited for generating interpretable summaries, whereas BERTopic better preserves structural relationships in the data. The choice of method should therefore depend on the specific goal of the analysis.

Despite these findings, there are also several limitations to consider. The data is often noisy and unstructured, which limits model performance. In addition, many tickets contain multiple issues, making it difficult for models to assign them to a single category. Evaluation metrics such as coherence and silhouette do not fully capture human interpretability, meaning that assessment is partly subjective.

From a practical perspective,the findings demonstrate the practical value of text analytics for understanding customer feedback. By identifying recurring patterns in support tickets, it becomes easier to pinpoint the types of issues that affect users most frequently. 
